% eksperymentalne tłumaczenie na włoski

%

\subsection{Prosta Wallace'a-Simsona} % lang-pl
\subsection{Retta di Wallace-Simson} % lang-it

\begin{proposition}
[twierdzenie Wallace'a] % lang-pl
[teorema di Wallace] % lang-it
	Dany jest trójkąt $ABC$ oraz punkt $P$. % lang-pl
	Niech $P_a$, $P_b$, $P_c$ oznacza rzuty punktu $P$ na proste $BC$, $AC$, $AB$. % lang-pl
	Wtedy następujące warunki są równoważne: punkt $P$ leży na okręgu opisanym na $\triangle ABC$; punkty $P_a$, $P_b$, $P_c$ są współliniowe. % lang-pl
	Siano dati un triangolo $ABC$ e un punto $P$. % lang-it
	Siano $P_a$, $P_b$, $P_c$ le proiezioni del punto $P$ sulle rette $BC$, $AC$, $AB$. % lang-it
	Allora le seguenti condizioni sono equivalenti: il punto $P$ appartiene alla circonferenza circoscritta a $\triangle ABC$; i punti $P_a$, $P_b$, $P_c$ sono allineati. % lang-it
\end{proposition}

Kiedy punkt $P$ leży na wspomnianym okręgu, prostą przez punkty $P_a$, $P_b$ oraz $P_c$ nazywa się prostą Simsona; na cześć Roberta Simsona, szkockiego matematyka żyjącego w latach 1687-1768; chociaż pierwszą pracę na jej temat opublikuje Wallace w 1799 roku. % lang-pl
Quando il punto $P$ appartiene alla circonferenza suddetta, la retta passante per i punti $P_a$, $P_b$ e $P_c$ è detta retta di Simson, in onore di Robert Simson, matematico scozzese vissuto tra il 1687 e il 1768, sebbene il primo lavoro sull'argomento sia stato pubblicato da Wallace nel 1799. % lang-it
\index[persons]{Wallace, ?}

Być może właściwsze byłoby pisać ,,prosta Wallace'a-Simsona''. % lang-pl
Forse sarebbe più corretto parlare di ``retta di Wallace-Simson''. % lang-it

Współliniowość spodków punktu $P$ na boki trójkąta jest twierdzeniem z dowodem u Zetela \cite[s. 55]{zetel_2020}; ćwiczeniem u Audina \cite[s. 104]{audin_2003}, Hartshorne'a \cite[s. 61]{hartshorne2000} i Neugebauera \cite[s. 25]{neugebauer_2018}. % lang-pl
L'allineamento dei piedi delle perpendicolari dal punto $P$ ai lati del triangolo è presentato come teorema con dimostrazione da Zetel \cite[p. 55]{zetel_2020}; come esercizio da Audin \cite[p. 104]{audin_2003}, Hartshorne \cite[p. 61]{hartshorne2000} e Neugebauer \cite[p. 25]{neugebauer_2018}. % lang-it
% \todofoot{HARTSHORNE PRZECHODZI DO PUNKTU MIQUELA}
Znajdzie się też w \cite[s. 24]{komisarczyk_2024}. % lang-pl
Pojawi się wielokrotnie w Delcie: w artykułach Michała Kiezy ($\Delta_{11}^9$), Joanny Jaszuńskiej ($\Delta_{15}^{10}$), Dominika Burka i Tomasza Cieśli ($\Delta_{16}^{11}$, z twierdzeniem Steinera). % lang-pl
Compare anche in \cite[p. 24]{komisarczyk_2024}. % lang-it
Compare più volte anche su \emph{Delta}: negli articoli di Michał Kieza ($\Delta_{11}^9$), Joanna Jaszuńska ($\Delta_{15}^{10}$), Dominik Burek e Tomasz Cieśla ($\Delta_{16}^{11}$, con il teorema di Steiner). % lang-it

\begin{proposition}
[twierdzenie Steinera] % lang-pl
[teorema di Steiner] % lang-it
	Dany jest trójkąt $ABC$ oraz punkt $P$ leżący na okręgu opisanym na trójkącie. % lang-pl
	Wtedy obrazy punktu $P$ w symetrii względem boków trójkąta leżą na jednej prostej, która przechodzi przez ortocentrum trójkąta. % lang-pl
	Siano dati un triangolo $ABC$ e un punto $P$ appartenente alla sua circonferenza circoscritta. % lang-it
	Allora le immagini del punto $P$ rispetto alle simmetrie aventi come assi i lati del triangolo giacciono su una stessa retta, che passa per l'ortocentro del triangolo. % lang-it
\end{proposition}

% PROSTA STEINERA? Audin ta sama strona
% prosta Steinera -> Znajdzie się też w \cite[s. 24]{komisarczyk_2024}.
% http://users.math.uoc.gr/~pamfilos/eGallery/problems/SteinerLine.html

Z twierdzenia o prostej Wallace'a-Simsona wynika, jak u Zetela \cite[s. 57]{zetel_2020}: % lang-pl
Dal teorema della retta di Wallace-Simson segue, come in Zetel \cite[p. 57]{zetel_2020}: % lang-it

\begin{proposition}
[twierdzenie Salmona] % lang-pl
[teorema di Salmon] % lang-it
\index{twierdzenie!Salmona}% % lang-pl
	Dany jest okrąg oraz trzy jego różne cięciwy $PA$, $PB$, $PC$ takie, że przekrojem okręgów na średnicach $PA$, $PB$ (odpowiednio: $PB$, $PC$ i $PA$, $PC$) są punkty $P$, $M$ (odpowiednio: $P$, $K$ oraz $P$, $L$). % lang-pl
	Wtedy punkty $K$, $L$, $M$ są współliniowe. % lang-pl
\index{teorema!di Salmon}% % lang-it
	Sono dati un cerchio e tre sue corde distinte $PA$, $PB$, $PC$ tali che l'intersezione dei cerchi aventi per diametri $PA$, $PB$ (rispettivamente $PB$, $PC$ e $PA$, $PC$) sia costituita dai punti $P$, $M$ (rispettivamente $P$, $K$ e $P$, $L$). % lang-it
	Allora i punti $K$, $L$, $M$ sono allineati. % lang-it
\end{proposition}

%